\textbf{Phần 1. Trắc nghiệm lựa chọn}
\setcounterpageref{ex}{0}

\begin{ex} %[3D1V45]
Số gồm $5$ nghìn, $8$ trăm, $0$ chục và $4$ đơn vị được viết là
\choice
{$5840$}
{$5084$}
{\True $5804$}
{$584$}
\loigiai{
Chữ số hàng nghìn là 5, hàng trăm là 8, hàng chục là 0, hàng đơn vị là 4. Số đó là $5804$.
\textbf{Đáp án của câu hỏi là C}
}
\end{ex}

\begin{ex} %[3D1V45]
Số liền sau của số $9999$ là
\choice
{\True $10\,000$}
{$9998$}
{$10\,001$}
{$9900$}
\loigiai{
Số liền sau của 9999 là: $9999 + 1 = 10\,000$.
\textbf{Đáp án của câu hỏi là A}
}
\end{ex}

\begin{ex} %[3D1V45]
Từ bốn chữ số $0; 3; 5; 8$, viết được số có bốn chữ số khác nhau lớn nhất là
\choice
{$8503$}
{$8350$}
{$5830$}
{\True $8530$}
\loigiai{
Chọn chữ số lớn nhất ở hàng cao nhất: hàng nghìn là 8, hàng trăm là 5, hàng chục là 3, hàng đơn vị là 0. Số đó là $8530$.
\textbf{Đáp án của câu hỏi là D}
}
\end{ex}

\begin{ex} %[3D1V45]
Số có bốn chữ số mà chữ số hàng nghìn gấp đôi chữ số hàng trăm, chữ số hàng chục gấp đôi chữ số hàng đơn vị và tổng các chữ số bằng 9 là
\choice
{$6310$}
{\True $4221$}
{$2142$}
{$4212$}
\loigiai{
Xét số 4221: nghìn 4 gấp đôi trăm 2; chục 2 gấp đôi đơn vị 1; tổng các chữ số là $4 + 2 + 2 + 1 = 9$ (thỏa mãn).
\textbf{Đáp án của câu hỏi là B}
}
\end{ex}

\begin{ex} %[3D1V45]
Trong số $7852$, giá trị của chữ số 8 lớn hơn giá trị của chữ số 5 bao nhiêu đơn vị?
\choice
{$700$}
{$300$}
{\True $750$}
{$800$}
\loigiai{
Chữ số 8 ở hàng trăm có giá trị 800. Chữ số 5 ở hàng chục có giá trị 50. Chênh lệch: $800 - 50 = 750$ đơn vị.
\textbf{Đáp án của câu hỏi là C}
}
\end{ex}

\begin{ex} %[3D1V45]
Số chẵn bé nhất có bốn chữ số khác nhau là
\choice
{\True $1024$}
{$1000$}
{$1023$}
{$1234$}
\loigiai{
Chữ số hàng nghìn là 1, hàng trăm là 0, hàng chục là 2, hàng đơn vị là chữ số chẵn bé nhất khác 0 và 2, tức là 4. Số đó là $1024$.
\textbf{Đáp án của câu hỏi là A}
}
\end{ex}

\begin{ex} %[3D1V45]
Có bao nhiêu số tròn trăm có bốn chữ số bé hơn $2000$?
\choice
{$9$ số}
{$8$ số}
{$11$ số}
{\True $10$ số}
\loigiai{
Các số tròn trăm từ 1000 đến 1900 gồm: 1000, 1100, 1200, ..., 1900. Có tất cả $(1900 - 1000) : 100 + 1 = 10$ số.
\textbf{Đáp án của câu hỏi là D}
}
\end{ex}

\begin{ex} %[3D1V45]
Một số có bốn chữ số có chữ số hàng đơn vị là 5. Nếu xóa bỏ chữ số 5 đó thì ta được số mới kém số ban đầu $2489$ đơn vị. Số ban đầu là
\choice
{$2545$}
{\True $2765$}
{$2655$}
{$2855$}
\loigiai{
Gọi số mới là $A$. Số ban đầu là $A \times 10 + 5$. Hiệu: $A \times 9 + 5 = 2489 \Rightarrow A \times 9 = 2484 \Rightarrow A = 276$. Số ban đầu là $2765$.
\textbf{Đáp án của câu hỏi là B}
}
\end{ex}

\begin{ex} %[3D1V45]
Số liền trước của số chẵn bé nhất có bốn chữ số là
\choice
{$1001$}
{$1000$}
{\True $999$}
{$998$}
\loigiai{
Số chẵn bé nhất có bốn chữ số là 1000. Số liền trước của nó là: $1000 - 1 = 999$.
\textbf{Đáp án của câu hỏi là C}
}
\end{ex}

\begin{ex} %[3D1V45]
Tìm một số biết rằng chữ số hàng nghìn là số lớn nhất có một chữ số, chữ số hàng đơn vị là số bé nhất có một chữ số, các chữ số còn lại đều bằng 3. Số đó là
\choice
{\True $9330$}
{$9303$}
{$9033$}
{$9333$}
\loigiai{
Hàng nghìn là 9, hàng đơn vị là 0, hàng trăm và hàng chục là 3. Số đó là $9330$.
\textbf{Đáp án của câu hỏi là A}
}
\end{ex}

\begin{ex} %[3D1V45]
Cho dãy số: $1000, 1500, 2000, 2500, \dots$ Số hạng thứ 8 của dãy số là
\choice
{$4000$}
{$5000$}
{$3500$}
{\True $4500$}
\loigiai{
Quy luật cộng thêm 500. Số hạng thứ 8 là: $1000 + 500 \times 7 = 1000 + 3500 = 4500$.
\textbf{Đáp án của câu hỏi là D}
}
\end{ex}

\begin{ex} %[3D1V45]
Số $10\,000$ gấp số $1000$ bao nhiêu lần?
\choice
{$100$ lần}
{\True $10$ lần}
{$5$ lần}
{$2$ lần}
\loigiai{
$10\,000 : 1000 = 10$ (lần).
\textbf{Đáp án của câu hỏi là B}
}
\end{ex}

\textbf{Phần 2. Trắc nghiệm đúng sai}
\setcounterpageref{ex}{0}

\begin{ex} %[3D1V45]
Xét tính đúng sai của các nhận định về các số đến 10 000:
\choiceTF
{\True Số 10 000 là số bé nhất có năm chữ số.}
{Số lớn nhất có bốn chữ số là 10 000.}
{\True Trong số 4725, chữ số 4 có giá trị là 4000.}
{Số 8080 đọc là tám trăm tám mươi.}
\loigiai{
\begin{itemize}
    \item Ý a) Đúng: Định nghĩa.
    \item Ý b) Sai: Sai vì số lớn nhất có 4 chữ số là 9999.
    \item Ý c) Đúng: Chữ số hàng nghìn.
    \item Ý d) Sai: Sai vì đọc là tám nghìn không trăm tám mươi.
\end{itemize}
\begin{itemize}
    \item a) \textbf{Đúng}.
    \item b) \textbf{Sai}.
    \item c) \textbf{Đúng}.
    \item d) \textbf{Sai}.
\end{itemize}
}
\end{ex}

\begin{ex} %[3D1V45]
Cho các chữ số $2; 0; 7; 9$:
\choiceTF
{\True Số lớn nhất có 4 chữ số khác nhau ghép từ các chữ số trên là 9720.}
{\True Số bé nhất có 4 chữ số khác nhau ghép từ các chữ số trên là 2079.}
{Có thể ghép được 24 số có 4 chữ số khác nhau từ 4 chữ số trên.}
{\True Hiệu giữa số lớn nhất và số bé nhất lập được là: $9720 - 2079 = 7641$.}
\loigiai{
\begin{itemize}
    \item Ý a) Đúng: Ghép theo thứ tự giảm dần.
    \item Ý b) Đúng: Chữ số hàng nghìn khác 0, ghép tăng dần.
    \item Ý c) Sai: Sai vì chữ số 0 không đứng đầu được nên chỉ có $3 \times 3 \times 2 \times 1 = 18$ số.
    \item Ý d) Đúng: Phép trừ chính xác.
\end{itemize}
\begin{itemize}
    \item a) \textbf{Đúng}.
    \item b) \textbf{Đúng}.
    \item c) \textbf{Sai}.
    \item d) \textbf{Đúng}.
\end{itemize}
}
\end{ex}

\textbf{Phần 3. Trả lời ngắn}
\setcounterpageref{ex}{0}

\begin{ex} %[3D1V45]
Tìm số liền trước của số 8500.
\shortans{8499}
\loigiai{
\textbf{Đáp số:} $8499$\\
$8500 - 1 = 8499$.
}
\end{ex}

\begin{ex} %[3D1V45]
Viết số bé nhất có bốn chữ số khác nhau.
\shortans{1023}
\loigiai{
\textbf{Đáp số:} $1023$\\
Số bé nhất có 4 chữ số khác nhau là 1023.
}
\end{ex}

\begin{ex} %[3D1V45]
Tính hiệu giữa số lớn nhất có 4 chữ số và số bé nhất có 4 chữ số.
\shortans{8999}
\loigiai{
\textbf{Đáp số:} $8999$\\
$9999 - 1000 = 8999$.
}
\end{ex}

\begin{ex} %[3D1V45]
Tìm giá trị của chữ số 6 trong số 6742.
\shortans{6000}
\loigiai{
\textbf{Đáp số:} $6000$\\
Chữ số 6 ở hàng nghìn có giá trị là 6000.
}
\end{ex}

\begin{ex} %[3D1V45]
Có bao nhiêu số có 4 chữ số đều có chữ số tận cùng là 0 từ 1000 đến 1090?
\shortans{10}
\loigiai{
\textbf{Đáp số:} $10$\\
Các số: 1000, 1010, 1020, 1030, 1040, 1050, 1060, 1070, 1080, 1090. Có 10 số.
}
\end{ex}

\begin{ex} %[3D1V45]
Một số có 4 chữ số có tổng các chữ số bằng 2. Hỏi có bao nhiêu số như vậy?
\shortans{4}
\loigiai{
\textbf{Đáp số:} $4$\\
Các số thỏa mãn: 2000, 1100, 1010, 1001. Có 4 số.
}
\end{ex}
